              13th International Mathematics Competition for University Students
                                   Odessa, July 20-26, 2006
                                          First Day

Problem 1. Let f : R → R be a real function. Prove or disprove each of the following statements.
    (a) If f is continuous and range(f ) = R then f is monotonic.
    (b) If f is monotonic and range(f ) = R then f is continuous.
    (c) If f is monotonic and f is continuous then range(f ) = R.
(20 points)
Solution. (a) False. Consider function f (x) = x3 − x. It is continuous, range(f ) = R but, for example,
f (0) = 0, f ( 12 ) = − 83 and f (1) = 0, therefore f (0) > f ( 21 ), f ( 12 ) < f (1) and f is not monotonic.
    (b) True. Assume first that f is non-decreasing. For an arbitrary number a, the limits lim f and
                                                                                                               a−
lim f exist and lim f ≤ lim f . If the two limits are equal, the function is continuous at a. Otherwise,
a+                     a−    a+
if lim f = b < lim f = c, we have f (x) ≤ b for all x < a and f (x) ≥ c for all x > a; therefore
     a−                a+
range(f ) ⊂ (−∞, b) ∪ (c, ∞) ∪ {f (a)} cannot be the complete R.
   For non-increasing f the same can be applied writing reverse relations or g(x) = −f (x).
   (c) False. The function g(x) = arctan x is monotonic and continuous, but range(g) = (−π/2, π/2) 6= R.

Problem 2. Find the number of positive integers x satisfying the following two conditions:
     1. x < 102006 ;
     2. x2 − x is divisible by 102006 .
(20 points)
Solution 1. Let Sk = 0 < x < 10k x2 − x is divisible by 10k and s (k) = |Sk | , k ≥ 1. Let x =
                               

ak+1 ak . . . a1 be the decimal writing of an integer x ∈ Sk+1 , k ≥ 1. Then obviously y = ak . . . a1 ∈ Sk . Now,
                                                                                                               2
let y = ak . . .a1 ∈ Sk be fixed. Considering ak+1 as a variable digit, we have x2 − x = ak+1 10k + y −
 ak+1 10k + y = (y 2 − y) + ak+1 10k (2y − 1) + a2k+1 102k . Since y 2 − y = 10k z for an iteger z, it follows that
x2 − x is divisible by 10k+1 if and only if z + ak+1 (2y − 1) ≡ 0 (mod 10). Since y ≡ 3 (mod 10) is obviously
impossible, the congruence has exactly one solution. Hence we obtain a one-to-one correspondence between
the sets Sk+1 and Sk for every k ≥ 1. Therefore s (2006) = s (1) = 3, because S1 = {1, 5, 6} .
Solution 2. Since x2 − x = x(x − 1) and the numbers x and x − 1 are relatively prime, one of them must
be divisible by 22006 and one of them (may be the same) must be divisible by 52006 . Therefore, x must
satisfy the following two conditions:

                                              x ≡ 0 or 1     (mod 22006 );

                                              x ≡ 0 or 1     (mod 52006 ).
Altogether we have 4 cases. The Chinese remainder theorem yields that in each case there is a unique
solution among the numbers 0, 1, . . . , 102006 − 1. These four numbers are different because each two gives
different residues modulo 22006 or 52006 . Moreover, one of the numbers is 0 which is not allowed.
    Therefore there exist 3 solutions.

Problem 3. Let A be an n × n-matrix with integer entries and b1 , . . . , bk be integers satisfying det A =
b1 · . . . · bk . Prove that there exist n × n-matrices B1 , . . . , Bk with integer entries such that A = B1 · . . . · Bk
and det Bi = bi for all i = 1, . . . , k.
(20 points)
Solution. By induction, it is enough to consider the case m = 2. Furthermore, we can multiply A with
any integral matrix with determinant 1 from the right or from the left, without changing the problem.
Hence we can assume A to be upper triangular.

                                                            1
Lemma. Let A be an integral upper triangular matrix, and let b, c be integers satisfying det A = bc. Then
there exist integral upper triangular matrices B, C such that det B = b, det C = c, A = BC.
Proof. The proof is done by induction on n, the case n = 1 being obvious. Assume the statement is true
for n − 1. Let A, b, c as in the statement of the lemma. Define Bnn to be the greatest common divisor of b
                             Ann
and Ann , and put Cnn = B     nn
                                 . Since Ann divides bc, Cnn divides Bbnn c, which divides c. Hence Cnn divides
c. Therefore, b0 = Bbnn and c0 = Cnn c
                                        are integers. Define A0 to be the upper-left (n − 1) × (n − 1)-submatrix
of A; then det A0 = b0 c0 . By induction we can find the upper-left (n − 1) × (n − 1)-part of B and C in such
a way that det B = b, det C = c and A = BC holds on the upper-left (n − 1) × (n − 1)-submatrix of A. It
remains to define Bi,n and Ci,n such that A = BC also holds for the (i, n)-th entry for all i < n.
    First we check that Bii and Cnn are relatively prime for all i < n. Since Bii divides b0 , it is certainly
enough to prove that b0 and Cnn are relatively prime, i.e.
                                                                      
                                                  b          Ann
                                       gcd               ,               = 1,
                                             gcd(b, Ann ) gcd(b, Ann )

which is obvious. Now we define Bj,n and Cj,n inductively: Suppose we have defined Bi,n and Ci,n for all
i = j + 1, j + 2, . . . , n − 1. Then Bj,n and Cj,n have to satisfy

                                 Aj,n = Bj,j Cj,n + Bj,j+1 Cj+1,n + · · · + Bj,n Cn,n

Since Bj,j and Cn,n are relatively prime, we can choose integers Cj,n and Bj,n such that this equation is
satisfied. Doing this step by step for all j = n − 1, n − 2, . . . , 1, we finally get B and C such that A = BC.
2

Problem 4. Let f be a rational function (i.e. the quotient of two real polynomials) and suppose that
f (n) is an integer for infinitely many integers n. Prove that f is a polynomial.
(20 points)
Solution. Let S be an infinite set of integers such that rational function f (x) is integral for all x ∈ S.
    Suppose that f (x) = p(x)/q(x) where p is a polynomial of degree k and q is a polynomial of degree n.
Then p, q are solutions to the simultaneous equations p(x) = q(x)f (x) for all x ∈ S that are not roots of
q. These are linear simultaneous equations in the coefficients of p, q with rational coefficients. Since they
have a solution, they have a rational solution.
    Thus there are polynomials p0 , q 0 with rational coefficients such that p0 (x) = q 0 (x)f (x) for all x ∈ S that
are not roots of q. Multiplying this with the previous equation, we see that p0 (x)q(x)f (x) = p(x)q 0 (x)f (x)
for all x ∈ S that are not roots of q. If x is not a root of p or q, then f (x) 6= 0, and hence p0 (x)q(x) =
p(x)q 0 (x) for all x ∈ S except for finitely many roots of p and q. Thus the two polynomials p0 q and pq 0
are equal for infinitely many choices of value. Thus p0 (x)q(x) = p(x)q 0 (x). Dividing by q(x)q 0 (x), we see
that p0 (x)/q 0 (x) = p(x)/q(x) = f (x). Thus f (x) can be written as the quotient of two polynomials with
rational coefficients. Multiplying up by some integer, it can be written as the quotient of two polynomials
with integer coefficients.
    Suppose f (x) = p00 (x)/q 00 (x) where p00 and q 00 both have integer coefficients. Then by Euler’s division
algorithm for polynomials, there exist polynomials s and r, both of which have rational coefficients such
that p00 (x) = q 00 (x)s(x) + r(x) and the degree of r is less than the degree of q 00 . Dividing by q 00 (x), we get
that f (x) = s(x) + r(x)/q 00 (x). Now there exists an integer N such that N s(x) has integral coefficients.
Then N f (x) − N s(x) is an integer for all x ∈ S. However, this is equal to the rational function N r/q 00 ,
which has a higher degree denominator than numerator, so tends to 0 as x tends to ∞. Thus for all
sufficiently large x ∈ S, N f (x) − N s(x) = 0 and hence r(x) = 0. Thus r has infinitely many roots, and is
0. Thus f (x) = s(x), so f is a polynomial.

Problem 5. Let a, b, c, d, e > 0 be real numbers such that a2 + b2 + c2 = d2 + e2 and a4 + b4 + c4 = d4 + e4 .
Compare the numbers a3 + b3 + c3 and d3 + e3 .
(20 points)


                                                          2
Solution. Without loss of generality a ≥ b ≥ c, d ≥ e. Let c2 = e2 + ∆, ∆ ∈ R. Then d2 = a2 + b2 + ∆
and the second equation implies
                                                                                               a2 b 2
                                a4 + b4 + (e2 + ∆)2 = (a2 + b2 + ∆)2 + e4 , ∆ = − a2 +b          2 −e2 .
         2      2    2        2 2      2       2       1 2     2    1 2     2
(Here a + b − e ≥ 3 (a + b + c ) − 2 (d + e ) = 6 (d + e ) > 0.)
                       a2 b 2       (a2 −e2 )(e2 −b2 )
Since c2 = e2 − a2 +b    2 −e2 =       a2 +b2 −e2
                                                        > 0 then a > e > b.
              2      2        2      a2 b 2         2
Therefore d = a + b − a2 +b2 −e2 < a and a > d ≥ e > b ≥ c.
     Consider a function f (x) = ax + bx + cx − dx − ex , x ∈ R. We shall prove that f (x) has only two
zeroes x = 2 and x = 4 and changes the sign at these points. Suppose the contrary. Then Rolle’s
theorem implies that f 0 (x) has at least two distinct zeroes. Without loss of generality a = 1. Then
f 0 (x) = ln b · bx + ln c · cx − ln d · dx − ln e · ex , x ∈ R. If f 0 (x1 ) = f 0 (x2 ) = 0, x1 < x2 , then
                                    ln b · bxi + ln c · cxi = ln d · dxi + ln e · exi , i = 1, 2,
but since 1 > d ≥ e > b ≥ c we have
                 (− ln b) · bx2 + (− ln c) · cx2     x2 −x1     x2 −x1   (− ln d) · dx2 + (− ln e) · ex2
                                                 ≤ b        < e        ≤                                 ,
                 (− ln b) · bx1 + (− ln c) · cx1                         (− ln d) · dx1 + (− ln e) · ex1
a contradiction. Therefore f (x) has a constantS sign at each of the intervals (−∞, 2), (2, 4) and (4, ∞).
Since f (0) = 1 then f (x) > 0, x ∈ (−∞, 2) (4, ∞) and f (x) < 0, x ∈ (2, 4). In particular, f (3) =
a3 + b3 + c3 − d3 − e3 < 0.
Problem 6. Find all sequences a0 , a1 , . . . , an of real numbers where n ≥ 1 and an 6= 0, for which the
following statement is true:
    If f : R → R is an n times differentiable function and x0 < x1 < . . . < xn are real numbers such that
f (x0 ) = f (x1 ) = . . . = f (xn ) = 0 then there exists an h ∈ (x0 , xn ) for which
                                      a0 f (h) + a1 f 0 (h) + . . . + an f (n) (h) = 0.
(20 points)
Solution. Let A(x) = a0 + a1 x + . . . + an xn . We prove that sequence a0 , . . . , an satisfies the required
property if and only if all zeros of polynomial A(x) are real.
   (a) Assume that all roots of A(x) are real. Let us use the following notations. Let I be the identity
operator on R → R functions and D be differentiation operator. For an arbitrary polynomial P (x) =
p0 + p1 x + . . . + pn xn , write P (D) = p0 I + p1 D + p2 D 2 + . . . + pn D n . Then the statement can written as
(A(D)f )(ξ) = 0.
   First prove the statement for n = 1. Consider the function
                                                              a0
                                                                   x
                                                    g(x) = e a1 f (x).
Since g(x0 ) = g(x1 ) = 0, by Rolle’s theorem there exists a ξ ∈ (x0 , x1 ) for which
                                                                       a0
                                                                            ξ
                                 a 0 a0 ξ        a0
                                                    ξ        e a1
                           0
                          g (ξ) = e a1 f (ξ) + e a1 f 0 ξ) =      (a0 f (ξ) + a1 f 0 (ξ)) = 0.
                                 a1                           a1
      Now assume that n > 1 and the statement holds for n−1. Let A(x) = (x−c)B(x) where c is a real root
of polynomial A. By the n = 1 case, there exist y0 ∈ (x0 , x1 ), y1 ∈ (x1 , x2 ), . . . , yn−1 ∈ (xn−1 , xn ) such that
f 0 (yj ) − cf (yj ) = 0 for all j = 0, 1, . . . , n − 1. Now apply the induction hypothesis for polynomial B(x),
function g = f 0 −cf and points y0 , . . . , yn−1 . The hypothesis says that there exists a ξ ∈ (y0 , yn−1 ) ⊂ (x0 , xn )
such that
                               (B(D)g)(ξ) = (B(D)(D − cI)f )(ξ) = (A(D)f )(ξ) = 0.
    (b) Assume that u + vi is a complex root of polynomial A(x) such that v 6= 0. Consider the linear
differential equation an g (n) + . . . + a1 g 0 + g = 0. A solution of this equation is g1 (x) = eux sin vx which has
infinitely many zeros.
    Let k be the smallest index for which ak 6= 0. Choose a small ε > 0 and set f (x) = g1 (x) + εxk . If
ε is sufficiently small then g has the required number of roots but a0 f + a1 f 0 + . . . + an f (n) = ak ε 6= 0
everywhere.

                                                             3
